\paragraph{Endpoint shielding budgets and a finite collision test.}
\label{par:terminal-budget-audit}
The signed shielding identity is exact, but its proposed
integrated bound must not be confused with an independently
proved upper estimate. The distinction matters at both
interior collisions and the terminal heat time.

\begin{proposition}[The endpoint budget is an identity]
\label{prop:budget-is-gap}
Let $a(s)<b(s)$ be consecutive real simple zeros of
$H_{1/2-s}$ on $0\le s<S$ and write
$G(s)=(b(s)-a(s))^2$ and
$\mathcal D(s)=2-G(s)\mathcal S(s)$.
For every $0\le\sigma<S$,
\begin{equation}
4\int_0^\sigma\mathcal D(s)\,ds
=G(0)-G(\sigma).
\label{eq:terminal-identity}
\end{equation}
Consequently
\begin{equation}
4\int_0^\sigma\mathcal D(s)\,ds<G(0)
\quad\Longleftrightarrow\quad G(\sigma)>0.
\label{eq:strict-budget-equivalent}
\end{equation}
A strict inequality through $\sigma=1/2$, if the
zero paths are defined there, is stronger than the
Riemann hypothesis: it also requires simplicity of
the corresponding critical-line zeros.
\end{proposition}
\begin{proof}
Equation~\eqref{eq:shieldgap} is $G'=-4\mathcal D$;
the fundamental theorem of calculus gives
\eqref{eq:terminal-identity}, and rearrangement gives
\eqref{eq:strict-budget-equivalent}.
RH places zeros on the critical line but does not
exclude multiple real zeros. A multiple endpoint
zero can arise from two distinct simple branches
whose squared gap tends to zero as $s\uparrow1/2$.
Thus demanding a strictly positive terminal gap
adds the unproved simplicity condition.
\end{proof}

The correct \emph{interior} sufficient target is
not the terminal strict budget itself but an
independent, pair-dependent envelope
$\mathcal B_i(\sigma)$ whose proof does not assume
survival of that pair to $\sigma$, together with
\begin{equation}
\boxed{
4\mathcal B_i(\sigma)\le
G_i(0)-\eta_i(\sigma),\qquad
\eta_i(\sigma)>0\quad
(0\le\sigma<1/2).
}
\label{eq:noncircular-shielding}
\end{equation}
If an a priori estimate also gives
$\int_0^\sigma\mathcal D_i\le\mathcal B_i(\sigma)$
up to any hypothetical first collision, then
\eqref{eq:terminal-identity} precludes that
finite-gap collision at any interior time.
The margin $\eta_i(\sigma)$ is allowed to
approach zero at $1/2$ and must be scaled
to the initial gap; no universal absolute
positive margin across all high zeros is
assumed. The independent envelope,
valid for \emph{all genuine theta-flow pairs},
has not been constructed. A separate
completeness argument must also rule out
a loss of the global real-zero property
through zeros whose heights escape to
infinity as the heat parameter varies.

\begin{proposition}[An exact quartic heat-flow
collision]\label{prop:quartic-shielding-control}
Consider the polynomial heat evolution
\begin{equation}
p_s(z)=e^{s\partial_z^2}
\bigl[(z^2-1)(z^2-4)\bigr]
=z^4+(12s-5)z^2+12s^2-10s+4.
\label{eq:quartic-heat-model}
\end{equation}
Initially its zeros are $-2,-1,1,2$,
all simple and real. The two positive zeros
collide at
\begin{equation}
s_*=\frac{10-\sqrt{46}}{24}
\in\left(\frac18,\frac17\right)
\subset(0,\tfrac12).
\label{eq:quartic-collision}
\end{equation}
For their initial gap $g(0)=1$, the two
exterior zeros contribute
$\mathcal S(0)=\frac14$, so
$\mathcal D(0)=\frac74$ and $G'(0)=-7$.
The integrated signed shielding defect reaches
\begin{equation}
4\int_0^{s_*}\mathcal D(s)\,ds
=G(0)=1.
\label{eq:quartic-saturation}
\end{equation}
This polynomial is an adversarial
heat-flow control, \emph{not} the
Riemann theta transform.
\end{proposition}
\begin{proof}
Since $\partial_z^2(z^4-5z^2+4)=12z^2-10$
and $\partial_z^4(z^4-5z^2+4)=24$,
the heat series terminates and gives
\eqref{eq:quartic-heat-model}.
Writing $v=z^2$, its quadratic
discriminant is
\[
\Delta(s)=(5-12s)^2
-4(12s^2-10s+4)
=96s^2-80s+9.
\]
The smaller root of $\Delta=0$
is \eqref{eq:quartic-collision}; it is
strictly between $1/8$ and $1/7$
and the two $v$-roots remain positive
until the collision. With
$B(s)=12s^2-10s+4>0$, the squared gap
between the positive roots is
\[
G(s)=5-12s-2\sqrt{B(s)}.
\]
It starts at $G(0)=1$ and tends to
$0$ at $s_*$. At $s=0$, the left
exterior roots $-1,-2$ contribute
\[
\mathcal S(0)
=\frac1{(1+1)(2+1)}
+\frac1{(1+2)(2+2)}
=\frac16+\frac1{12}=\frac14.
\]
The zero-motion equation gives
$\mathcal D(0)=2-\frac14=\frac74$
and $G'(0)=-7$.
The fundamental theorem of calculus
applied up to the collision gives
\eqref{eq:quartic-saturation}.
\end{proof}

The control shows that a positive exterior
shielding sum alone does not guarantee
continuation of simple real zeros:
what matters is its signed accumulated
deficit relative to the \emph{specific}
initial gap. Moreover
\eqref{eq:quartic-saturation} explains the
strictness required on open heat-time
intervals and why an endpoint equality
should remain permissible. The accompanying
exact-fraction regression tests $2{,}000$
quartic identities and $1{,}000$ independent
budget identities; a separate $70$-digit
diagnostic checks the collision time and
integral. See
\path{certificates/xi_signed_shielding_budget_audit.py}
and \path{XiKernel/SignedShieldingEndpoint.lean}.
The Lean declarations certify the algebraic
gates, not the whole theta zero-motion
ODE or a universal signed-defect estimate.
The full Riemann hypothesis remains open.
